\documentclass[conference]{IEEEtran}

\usepackage{amsmath,amssymb,bm}
\usepackage{booktabs}
\usepackage{graphicx}
\usepackage{cite}
\usepackage{xcolor}
\usepackage{url}
\usepackage{multirow}

\newcommand{\todo}[1]{\textcolor{red!70!black}{\textbf{[TODO:} #1\textbf{]}}}
\newcommand{\calD}{\mathcal{D}}
\newcommand{\feasSet}{\mathcal{F}}
\newcommand{\AURCstar}{\mathrm{AURC}_{\star}}
\newcommand{\Nfive}{N_{5\%}}
\newcommand{\distdyn}{d^{\mathrm{dyn}}}
\newcommand{\placeholderfigurewide}[1]{%
    \fbox{\parbox[c][4.5cm][c]{0.96\textwidth}{\centering\small
    \textbf{FIGURE PLACEHOLDER}\\[1.5mm]#1}}}

\title{Fleet-Informed Bayesian Optimization for Data-Efficient Controller Calibration}

\author{\IEEEauthorblockN{Jakob Weber, Markus Gurtner, Adrian Trachte, Andreas Kugi}
\IEEEauthorblockA{\textit{Draft author list -- affiliations to be finalized}}}

\begin{document}
\maketitle

\begin{abstract}
Controller calibration is commonly repeated for every new system variant, although mature fleets may already contain substantial information about which controller-design parameters perform well. This paper studies whether historical fleet calibration experience can reduce the number of new closed-loop experiments required for a newly arriving client. We consider LQI yaw-rate tracking for a heterogeneous vehicle fleet and use Bayesian optimization (BO) to tune a three-dimensional controller-weight vector. Historical scalar calibration observations from previously optimized clients warm-start the target client's Gaussian-process surrogate, while every candidate controller is synthesized from the target client's own identified model and evaluated only on its local nonlinear plant. A pre-registered confirmatory campaign with 20 previously untouched fleet seeds and 120 held-out targets shows that mature-fleet warm start reduces the mean number of target-local experiments required to reach within 5\% of the client-specific oracle from 6.63 to 3.07, a 54\% reduction, while strongly reducing variability. Dynamics-aware weighting of historical observations provides no measurable improvement over unweighted pooling. A separate landscape study explains this result: despite substantial plant heterogeneity, pairwise LQI-calibration landscapes remain highly rank-correlated, with Spearman correlation above 0.95 for every evaluated client pair. The results indicate that fleet-level controller-calibration knowledge can be broadly reusable even when the underlying plant dynamics differ substantially.
\end{abstract}

\begin{IEEEkeywords}
Bayesian optimization, controller calibration, transfer learning, LQI, heterogeneous systems, vehicle control
\end{IEEEkeywords}

\section{Introduction}
Controller calibration is an important but expensive part of commissioning modern dynamical systems. Even when the controller structure is fixed by model-based design, parameters such as LQR/LQI state and input penalties still determine the final compromise between tracking performance, control effort, and robustness. Bayesian optimization (BO) is attractive for this task because it can optimize low-dimensional controller parameters from a limited number of expensive zeroth-order evaluations~\cite{frazier2018tutorial,marco2016lqr,berkenkamp2016safe,schillinger2017safe,vonrohr2024local}.

A fleet setting changes the calibration problem fundamentally. Vehicles or machines are commissioned one after another, and earlier clients may already have accumulated closed-loop calibration experience. Repeating a cold-start optimization for every new client ignores this information. Collaborative and federated BO similarly exploit information distributed across multiple agents~\cite{dai2020fts,dai2021dpfts,chen2025collaborative}. For controller calibration, however, the central question is how transferable previous calibration observations remain when the underlying dynamical systems differ.

The intuitive answer is not obvious. Heterogeneous vehicle dynamics can change closed-loop behavior substantially, and uncertainty-aware dynamics estimates contain useful structure for clustered system identification and controller design~\cite{weber2026ifac}. This motivates dynamics-aware transfer. At the same time, the controller-calibration problem may be much more invariant than the plant itself: different models can induce similar rankings of the same low-dimensional LQI weight configurations, in which case broad reuse of fleet experience may be sufficient.

This paper studies that distinction in a mature-fleet setting. A new target joins a fleet whose previous clients have already completed BO-based controller calibration. Historical clients provide scalar calibration records in the common controller-weight space. The target starts with two local initialization experiments, fits a GP surrogate jointly to its own and historical observations, and sequentially evaluates new target-specific LQI controllers. Historical observations do not count toward the target's experimental budget, and no fixed feedback gain is transferred.

The contributions are:
\begin{itemize}
    \item We formulate \emph{mature-fleet controller calibration} as a transfer problem in which historical fleet calibration observations reduce the number of new closed-loop experiments required by a newly arriving client.
    \item In a pre-registered held-out campaign on a heterogeneous nonlinear vehicle fleet, fleet-informed warm start reduces the mean number of local evaluations needed to reach within 5\% of the client-specific oracle from 6.63 to 3.07 and strongly reduces the difficult-client tail.
    \item We characterize why broad transfer is possible. Although identified vehicle dynamics differ substantially, the LQI-calibration landscapes remain highly rank-correlated. Consistent with this observation, dynamics-aware weighting does not outperform simple unweighted pooling on fresh confirmatory data.
\end{itemize}

Figure~\ref{fig:concept} summarizes the commissioning workflow.

\begin{figure*}[t]
    \centering
    \placeholderfigurewide{\textbf{Fig. 1 -- Fleet-informed controller calibration concept.} Mature fleet with historical calibration records $\rightarrow$ newly arriving target $\rightarrow$ two target-local initialization experiments $\rightarrow$ fleet-informed GP/EI $\rightarrow$ target-local LQI synthesis and closed-loop evaluation.}
    \caption{Fleet-informed controller calibration. Historical calibration records from mature clients warm-start the BO surrogate of a newly arriving target. Only the target's new closed-loop evaluations count toward the commissioning budget; controller gains are synthesized and evaluated locally for the target.}
    \label{fig:concept}
\end{figure*}

